\section{Proper-subface flagged CRT surplus and endpoint residue cubes}
\label{sec:flagged-crt-residue-cube}

This section refines the shared-CRT boundary of
Section~\ref{sec:shared-crt-incidence-successor}.  The refinement has two
purposes.  First, it identifies and removes a false source of apparent
progress: if the surplus of a saturated block may be sent to an arbitrary
residual source, the resulting optimization is exactly the scalar endpoint
defect.  Second, it retains surplus only when a proper subpacket supplies an
independent prime-power congruence for the target source.  The resulting
finite boundary, denoted $B_{\rm FCRT}$, obeys the same once-charge height
inequality, can be strictly smaller than $B_{\rm SCRT}$, and need not equal
the scalar defect.  The associated uniform estimate is left as an explicit
open premise throughout.

\subsection{Endpoint data}

Let $P=(a,b,c)$ satisfy
\[
 a,b,c\in\mathbb Z_{>0},\qquad a+b=c,\qquad \gcd(a,b)=1.
\]
Pairwise coprimality follows.  Put
\[
 I=\{p:p\mid c,\ e_p:=v_p(c)\ge 2\},\qquad
 J=\{q:q\mid ab\},
\]
and assign the logarithmic source and sink weights
\[
 x_p=(e_p-1)\log p,\qquad y_q=\log q,
 \qquad X=\sum_{p\in I}x_p,
 \qquad Y=\sum_{q\in J}y_q.
\tag{FCRT-1}\label{eq:fcrt-endpoint-data}
\]
For $S\subseteq I$ and $T\subseteq J$, define
\[
 \begin{aligned}
 m(S)&=\prod_{p\in S}p^{e_p},
 &x(S)&=\sum_{p\in S}x_p,\\
 a_T&=\prod_{\substack{q\in T\\q\mid a}}q^{v_q(a)},
 &b_T&=\prod_{\substack{q\in T\\q\mid b}}q^{v_q(b)},\\
 y(T)&=\sum_{q\in T}y_q.
 \end{aligned}
\tag{FCRT-2}\label{eq:fcrt-packets}
\]
Empty products are one.  A pair $(S,T)$ of nonempty subsets is a
\emph{shared CRT block} if
\[
                 m(S)\mid a_T+b_T,
 \qquad           x(S)\le y(T).
\tag{FCRT-3}\label{eq:fcrt-block}
\]
The two clauses will be called compatibility and saturation, respectively.
Let $\mathcal B(P)$ be the set of all shared CRT blocks.  It is a finite
subset of
\[
 \bigl(\mathcal P(I)\setminus\{\varnothing\}\bigr)
 \times
 \bigl(\mathcal P(J)\setminus\{\varnothing\}\bigr).
\]
Prime factorization gives the exact identity
\[
 h(P)-r(P)
 =\log c-\log\rad(abc)=X-Y,
\tag{FCRT-4}\label{eq:fcrt-endpoint-identity}
\]
and hence the scalar lower bound is
\[
                         \Delta(P)=(X-Y)_+.
\]

\subsection{The exact collapse of free-target reuse}

It is useful to dispose of an unrestricted surplus rule before imposing the
arithmetic flag.  Let $K\subseteq\mathcal B(P)$ be a selected block family
such that the sets $S_\nu$ are pairwise disjoint and the sets $T_\nu$ are
pairwise disjoint, where $\nu=(S_\nu,T_\nu)\in K$.  In particular, the
blocks are distinct and $K$ is finite.
Write
\[
 I_0=I\setminus\bigcup_{\nu\in K}S_\nu,
 \qquad
 J_0=J\setminus\bigcup_{\nu\in K}T_\nu,
 \qquad
 \sigma_\nu=y(T_\nu)-x(S_\nu)\ge0.
\tag{FCRT-5}\label{eq:fcrt-free-residual-sets}
\]
A free-target configuration consists further of two total maps into pointed
finite sets,
\[
 o:J_0\longrightarrow I_0\sqcup\{\bot\},
 \qquad
 t:K\longrightarrow I_0\sqcup\{\bot\}.
\tag{FCRT-6}\label{eq:fcrt-free-owners}
\]
Here $o(q)=p$ assigns the residual sink $q$ exclusively to $p$, while
$t(\nu)=p$ sends the whole unsplit surplus $\sigma_\nu$ to $p$; the value
$\bot$ means unused.  Several independently charged objects may have the
same target.  Define
\[
 C_p=\sum_{\substack{q\in J_0\\o(q)=p}}y_q
     +\sum_{\substack{\nu\in K\\t(\nu)=p}}\sigma_\nu,
 \qquad
 B_{\rm free}(\mathcal C)=\sum_{p\in I_0}(x_p-C_p)_+.
\tag{FCRT-7}\label{eq:fcrt-free-boundary}
\]

\begin{proposition}[Free-target collapse]
\label{prop:fcrt-free-collapse}
For every positive primitive $P$,
\[
                 \min_{\mathcal C}B_{\rm free}(\mathcal C)
                 =(X-Y)_+.
\tag{FCRT-8}\label{eq:fcrt-free-collapse}
\]
\end{proposition}

\begin{proof}
Every sink weight is charged at most once.  A selected block consumes
$x(S_\nu)$ and leaves at most $\sigma_\nu$ for reuse, whereas an unselected
sink contributes at most its own weight.  Therefore
\[
                       X\le Y+B_{\rm free}(\mathcal C)
\]
for every configuration, proving the lower bound.

If $I=\varnothing$, the empty configuration realizes zero.  Suppose
$I\ne\varnothing$ and $X\le Y$.  Then $J\ne\varnothing$: otherwise
$a=b=1$, $c=2$, and $I$ would be empty.  The full pair $(I,J)$ is compatible
because $m(I)\mid c=a+b$, and it is saturated because $X\le Y$.  Selecting
it realizes zero.

It remains to consider $X>Y$.  Choose an inclusion-maximal subset
$S\subset I$ with $x(S)\le Y$.  It exists and is proper.  If
$S\ne\varnothing$, select the full-sink block $(S,J)$, which is compatible
and saturated.  For every $p\in I\setminus S$, maximality gives
\[
                         x_p>Y-x(S)=:\sigma.
\tag{FCRT-9}\label{eq:fcrt-maximality}
\]
Send the unique surplus token $\sigma$ to one such $p$.  Since no residual
sinks remain, all residual terms are positive and their sum is
\[
 x_p-\sigma+\sum_{u\in I\setminus(S\cup\{p\})}x_u=X-Y.
\]
If $S=\varnothing$, maximality says $x_p>Y$ for every $p\in I$.
Select no block and assign every residual sink to any one source.  The
boundary is again $X-Y$.  This proves equality in all cases.
\end{proof}

Proposition~\ref{prop:fcrt-free-collapse} rules out only this exact
free-target child model.  It supplies neither a counterexample to the
shared-CRT route nor a reason to remove arithmetic restrictions from a
surplus token.

\subsection{A fully typed proper-subface configuration}

Fix a finite block family as in \eqref{eq:fcrt-free-residual-sets}.  For each
$\nu\in K$, define its finite flag set
\[
 \mathcal F_\nu=
 \left\{(p,U)\ \middle|\
 \begin{array}{l}
 p\in I_0,\quad U\subseteq T_\nu,\quad
 \varnothing\ne U\ne T_\nu,\\[2pt]
 p^{e_p}\mid a_U+b_U
 \end{array}
 \right\}.
\tag{FCRT-10}\label{eq:fcrt-flag-type}
\]
Thus $U$ is a nonempty proper face of the block's sink packet, and the target
$p$ is residual, hence lies outside every selected block source set.  An
\emph{FCRT configuration} consists of the block family, a residual-sink
owner map
\[
                        o:J_0\longrightarrow I_0\sqcup\{\bot\},
\]
and, for every $\nu\in K$, a choice
\[
 f_\nu\in\mathcal F_\nu\sqcup\{\bot\}.
\tag{FCRT-11}\label{eq:fcrt-flag-choice}
\]
If $f_\nu=\bot$, put $\rho_\nu=0$.  If
$f_\nu=(p_\nu,U_\nu)$, put
\[
 \rho_\nu=\min\{\sigma_\nu,y(U_\nu)\},
 \qquad \sigma_\nu=y(T_\nu)-x(S_\nu).
\tag{FCRT-12}\label{eq:fcrt-reuse-cap}
\]
The token of mass $\rho_\nu$ is indivisible, has the single target $p_\nu$,
and travels one hop.  The amount $\sigma_\nu-\rho_\nu$ is discarded.  Credit
left after a target is filled is also discarded and cannot emit a new token.
Different residual sinks and different block tokens may meet at the same
source; this does not repeat the charge of any originating sink.

For $p\in I_0$, set
\[
 \begin{split}
 M_p&=\sum_{\substack{q\in J_0\\o(q)=p}}y_q,\\
 Z_p&=\sum_{\nu\in K}z_{\nu,p},
 \qquad
 z_{\nu,p}=
 \begin{cases}
  \rho_\nu,&f_\nu=(p,U)\text{ for some }U,\\
  0,&\text{otherwise},
 \end{cases}\\
 R_p&=(x_p-M_p-Z_p)_+,
 \end{split}
 \qquad
 B_{\rm FCRT}(\mathcal C)=\sum_{p\in I_0}R_p.
\tag{FCRT-13}\label{eq:fcrt-boundary}
\]
Every set and every choice space in
\eqref{eq:fcrt-flag-type}--\eqref{eq:fcrt-flag-choice} is finite.  The empty
block family with all residual sinks unused is admissible, so the pointwise
minimum
\[
 B_{\rm FCRT}(P)=\min_{\mathcal C\in\operatorname{Conf}_{\rm FCRT}(P)}
                         B_{\rm FCRT}(\mathcal C)
\tag{FCRT-14}\label{eq:fcrt-optimum}
\]
exists.

\begin{proposition}[Once-charge mass bridge]
\label{prop:fcrt-once-charge}
Every FCRT configuration satisfies
\[
                         X\le Y+B_{\rm FCRT}(\mathcal C)
\tag{FCRT-15}\label{eq:fcrt-mass-bridge}
\]
and consequently
\[
                         h(P)\le r(P)+B_{\rm FCRT}(\mathcal C).
\tag{FCRT-16}\label{eq:fcrt-height-bridge}
\]
\end{proposition}

\begin{proof}
The definition of the positive part gives
$x_p\le M_p+Z_p+R_p$.  Exclusive ownership of residual sinks gives
\[
                       \sum_{p\in I_0}M_p\le y(J_0).
\]
Each block emits at most one token, and
$\rho_\nu\le\sigma_\nu$, so
\[
                       \sum_{p\in I_0}Z_p
                       \le\sum_{\nu\in K}\sigma_\nu.
\]
Pairwise source disjointness, pairwise sink disjointness, and
\eqref{eq:fcrt-free-residual-sets} now yield
\[
 \begin{aligned}
 X
 &=\sum_{\nu\in K}x(S_\nu)+\sum_{p\in I_0}x_p\\
 &\le\sum_{\nu\in K}x(S_\nu)+y(J_0)
       +\sum_{\nu\in K}\bigl(y(T_\nu)-x(S_\nu)\bigr)
       +B_{\rm FCRT}(\mathcal C)\\
 &=Y+B_{\rm FCRT}(\mathcal C).
 \end{aligned}
\]
Combining this with \eqref{eq:fcrt-endpoint-identity} proves the height
inequality.
\end{proof}

The congruence and proper-face clauses do not enter the last accounting
calculation.  Their role is to restrict which once-charged configurations
exist; deleting them recovers the scalar collapse above.

\begin{corollary}[Pointwise sandwich]
\label{cor:fcrt-sandwich}
For every positive primitive point,
\[
 \Delta(P)\le B_{\rm FCRT}(P)
             \le B_{\rm SCRT}(P)\le B_{\rm PBT}(P).
\tag{FCRT-17}\label{eq:fcrt-sandwich}
\]
\end{corollary}

\begin{proof}
Every boundary is nonnegative, and
Proposition~\ref{prop:fcrt-once-charge} gives
$B_{\rm FCRT}(\mathcal C)\ge X-Y$ configuration by configuration.  Taking
the finite minimum therefore gives
$B_{\rm FCRT}(P)\ge\max\{X-Y,0\}=\Delta(P)$.  Giving every selected block
the choice $\bot$ embeds the SCRT configurations into the FCRT
configurations.  Taking no shared block recovers the exclusive PBT problem.
\end{proof}

\begin{definition}[Uniform proper-subface gate]
\label{def:fcrt-uniform-gate}
The statement \textup{FCRT-1} is
\[
 \forall\varepsilon>0\ \exists C_\varepsilon\in\mathbb R\
 \forall P\ \exists\mathcal C\in\operatorname{Conf}_{\rm FCRT}(P):
 \quad B_{\rm FCRT}(\mathcal C)
 \le\varepsilon r(P)+C_\varepsilon,
\tag{FCRT-18}\label{eq:fcrt-uniform-gate}
\]
where $P$ ranges over all positive primitive triples.
\end{definition}

\begin{theorem}[Conditional implication to $abc$]
\label{thm:fcrt-implies-abc}
If \textup{FCRT-1} holds, then the standard logarithmic $abc$ conjecture
holds.
\end{theorem}

\begin{proof}
Insert the configuration supplied by \eqref{eq:fcrt-uniform-gate} into
\eqref{eq:fcrt-height-bridge}:
\[
 h(P)\le r(P)+B_{\rm FCRT}(\mathcal C)
       \le(1+\varepsilon)r(P)+C_\varepsilon.
\]
The quantifier order is unchanged.
\end{proof}

No instance, inhabitant, or proof of the uniform gate
\eqref{eq:fcrt-uniform-gate} is supplied here.  The theorem is a conditional
reduction, not an unconditional proof of $abc$.

\subsection{Three exact separating points}

The following proposition records complete finite optimizations, rather than
merely exhibiting upper bounds.

\begin{proposition}[Strict improvement, noncollapse, and an active cap]
\label{prop:fcrt-three-witnesses}
The three positive primitive points below have the indicated exact
boundaries:
\[
\begin{array}{c|c|c|c|c}
 P&e^{\Delta(P)}&e^{B_{\rm FCRT}(P)}
  &e^{B_{\rm SCRT}(P)}&e^{B_{\rm PBT}(P)}\\ \hline
 (1,675,676)&26/15&26/15&2&2\\
 (1,224,225)&15/14&3/2&3/2&3/2\\
 (1,65024,65025)&255/254&3/2&3&15/2.
\end{array}
\tag{FCRT-19}\label{eq:fcrt-three-witness-table}
\]
In particular, the last row makes every inequality in
\eqref{eq:fcrt-sandwich} strict.
\end{proposition}

\begin{proof}
For $P=(1,675,676)$,
\[
 675=3^3 5^2,\qquad676=2^2 13^2,\qquad
 e^X=26,\quad e^Y=15.
\]
The proper packet $\{3\}$ satisfies
$1+3^3=28\equiv0\pmod4$, while $1+5^2=26$ is divisible by neither $4$ nor
$13^2$.  The full packet certifies both source moduli.  The complete list of
saturated compatible blocks is
\[
 (\{2\},\{3\}),\qquad(\{2\},\{3,5\}),\qquad
 (\{13\},\{3,5\}).
\]
They share sinks, so at most one is selectable.  The four complete
assignments of the two residual sinks in the no-block case have boundary
factors
\[
                         13,\qquad\frac{13}{5},\qquad
                         \frac{13}{3},\qquad2,
\]
while the three one-block cases have best factors
$13/5$, $13$, and $2$, in the order displayed above.  Thus the exact SCRT
minimum is $\log2$; the no-block list also gives the PBT minimum.  In FCRT,
choose
$(S,T)=(\{13\},\{3,5\})$.  Its surplus is
$\log(15/13)$, and $U=\{3\}$ is a legal face for residual source $2$.
Because $15/13<3$, the surplus, rather than the face weight, is limiting:
the whole surplus is reused and the residual factor is
\[
                         \frac{2}{15/13}=\frac{26}{15}.
\]
The scalar lower bound in \eqref{eq:fcrt-sandwich} proves optimality.

For $P=(1,224,225)$,
\[
 224=2^5\cdot7,\qquad225=3^2 5^2,\qquad
 e^X=15,\quad e^Y=14.
\]
The two nonempty proper packets give $1+2^5=33$ and $1+7=8$; neither is
divisible by $9$ or $25$.  Hence the only nontrivial saturated blocks use
the full sink set with one singleton source, and neither block can emit a
flagged token.  The four complete PBT assignments give residual factors
\[
                         5,\quad3,\quad\frac32,\quad\frac52.
\]
The full-sink singleton blocks leave factors $5$ and $3$.  Thus all three
finite models have exact factor $3/2$, strictly above the scalar factor
$15/14$.

Finally,
\[
 65024=2^9\cdot127,\qquad
 65025=3^2\cdot5^2\cdot17^2,\qquad
 e^X=255,\quad e^Y=254.
\]
The face $\{2\}$ gives
$1+2^9=513=3^3\cdot19$, so it certifies source $3$ and neither source $5$
nor source $17$.  The face $\{127\}$ gives $128$ and certifies none of the
three source moduli.  Every nonempty proper source subset paired with the
full sink set is compatible and saturated; the full source set is not
saturated.  The only compatible proper-packet block is
$(\{3\},\{2\})$, and it is not saturated.  Consequently every selectable
block uses the full sink set and at most one can be selected.

In PBT, the nine complete assignments of the two sinks to the three sources
have residual factors among
\[
 \frac{15}{2},\quad15,\quad\frac{51}{2},\quad51,
 \quad\frac{85}{2},\quad85;
\]
the minimum is $15/2$.  In SCRT, the block
$(\{5,17\},\{2,127\})$ closes source factor $85$ and leaves factor $3$.
Every other legal full-sink block leaves at least
$5,15,17,51$, or $85$, proving exactness.  The same block has surplus factor
$254/85$, while the legal face $U=\{2\}$ has factor $2$ and targets the
remaining source $3$.  Here the proper-face cap is active:
\[
 e^\rho=\min\left\{\frac{254}{85},2\right\}=2,
 \qquad e^{B_{\rm FCRT}}=\frac32.
\]
If a block contains source $3$, it cannot target $3$; if it omits source
$3$, the token just described is the only possible flagged credit and is at
most a factor of $2$.  The chosen block closes the largest possible mass in
the other sources.  This proves the matching lower bound.  Hence
\[
                         \frac{255}{254}<\frac32<3<\frac{15}{2}.
\]
\end{proof}

The middle row disproves the pointwise identity
$B_{\rm FCRT}=\Delta$.  It does not disprove the quantified uniform gate,
which permits both an $\varepsilon r(P)$ term and an additive constant.

\subsection{The endpoint residue cube}

Fix a source face $S\subseteq I$ and define the finite abelian group
\[
                   G_S=\prod_{p\in S}
                   (\mathbb Z/p^{e_p}\mathbb Z)^\times.
\tag{FCRT-20}\label{eq:fcrt-residue-group}
\]
For every integer $n$ coprime to $m(S)$, let $\overline n_S\in G_S$ denote
its diagonal residue tuple.  Since $\gcd(ab,c)=1$, the following elements are
well-defined:
\[
 g_{S,q}=
 \begin{cases}
  \overline{q^{v_q(a)}}_S,&q\mid a,\\
  \overline{q^{v_q(b)}}_S^{-1},&q\mid b.
 \end{cases}
\tag{FCRT-21}\label{eq:fcrt-residue-generators}
\]
Define the packet-label map on the Boolean cube by
\[
 \Phi_S:\mathcal P(J)\longrightarrow G_S,
 \qquad \Phi_S(T)=\prod_{q\in T}g_{S,q}.
\tag{FCRT-22}\label{eq:fcrt-packet-label}
\]

\begin{proposition}[Fixed-fibre and complement identities]
\label{prop:fcrt-fixed-fibre}
For every $T\subseteq J$,
\[
 \Phi_S(T)=\overline{a_T}_S\,\overline{b_T}_S^{-1},
 \qquad \Phi_S(J)=\overline{-1}_S,
\tag{FCRT-23}\label{eq:fcrt-full-label}
\]
and
\[
 \begin{split}
 m(S)\mid a_T+b_T
 &\quad\Longleftrightarrow\quad
 \Phi_S(T)=\overline{-1}_S\\
 &\quad\Longleftrightarrow\quad
 \Phi_S(J\setminus T)=1.
 \end{split}
\tag{FCRT-24}\label{eq:fcrt-complement-fibre}
\]
If $S'\subseteq S$ and
$\pi_{S,S'}:G_S\to G_{S'}$ is coordinate projection, then
\[
                         \pi_{S,S'}\circ\Phi_S=\Phi_{S'}.
\tag{FCRT-25}\label{eq:fcrt-inverse-system}
\]
\end{proposition}

\begin{proof}
Separating the primes of $T$ between the two coprime arms gives the first
identity in \eqref{eq:fcrt-full-label}.  At every coordinate $p^{e_p}$,
$a+b\equiv0$ and $b$ is a unit, so $ab^{-1}\equiv-1$; this gives the second
identity.  Since $b_T$ is a unit modulo every $p^{e_p}$,
$a_T+b_T\equiv0$ is equivalent to
$\overline{a_T}_S\overline{b_T}_S^{-1}=\overline{-1}_S$.
Finally,
\[
 \Phi_S(T)\Phi_S(J\setminus T)=\Phi_S(J)=\overline{-1}_S,
\]
which proves the complement statement by cancellation.  The projection
identity follows coordinatewise from the definitions.
\end{proof}

Thus compatible packets form the $\overline{-1}_S$ fibre, and complementation
identifies this fibre with the identity fibre.  The groups and maps form an
inverse system over the source-face poset.  An FCRT flag is nevertheless a
cross-fibre incidence.  For a block source $S_\nu$ and residual target
$p\notin S_\nu$, it requires
\[
 T_\nu\in\Phi_{S_\nu}^{-1}(\overline{-1}_{S_\nu}),
 \qquad
 U_\nu\in\Phi_{\{p\}}^{-1}(\overline{-1}_{\{p\}}),
 \qquad U_\nu\subsetneq T_\nu.
\tag{FCRT-26}\label{eq:fcrt-cross-fibre}
\]
Both labels can be compared by projection from
$G_{S_\nu\cup\{p\}}$, but neither packet is required to lie in the
$\overline{-1}$ fibre for the union source.  Consequently
\eqref{eq:fcrt-inverse-system} does not convert the flag into a fibre
inclusion.

\subsection{Signed exchange and exact Fourier enumerators}

Define the residue-relation lattice
\[
 \Lambda_S=\left\{z\in\mathbb Z^J:
              \prod_{q\in J}g_{S,q}^{z_q}=1\right\}.
\tag{FCRT-27}\label{eq:fcrt-residue-lattice}
\]

\begin{proposition}[Same-fibre signed exchange]
\label{prop:fcrt-signed-exchange}
If $T,T'\subseteq J$ are compatible with the same source face $S$, then
\[
 \prod_{q\in T'\setminus T}g_{S,q}
 =\prod_{q\in T\setminus T'}g_{S,q},
\tag{FCRT-28}\label{eq:fcrt-signed-product}
\]
and
\[
 y(T')-y(T)=y(T'\setminus T)-y(T\setminus T').
\tag{FCRT-29}\label{eq:fcrt-signed-weight}
\]
If $T\ne T'$, then
$\mathbf 1_{T'}-\mathbf 1_T$ is a nonzero member of
$\Lambda_S\cap\{-1,0,1\}^J$.
\end{proposition}

\begin{proof}
Both packet products are $\overline{-1}_S$.  Factor each over
$T\cap T'$ and its disjoint remainder, then cancel the common factor.  This
proves \eqref{eq:fcrt-signed-product}; the analogous disjoint decompositions
of the weight sums prove \eqref{eq:fcrt-signed-weight}.  The final statement
is exactly the multiplicative identity written in exponent-vector form.
\end{proof}

In particular, if $2^{|J|}>|G_S|$, the pigeonhole principle gives distinct
packets with the same label, and hence a nonzero
$\{-1,0,1\}$ relation in $\Lambda_S$.  This conclusion alone controls
neither the sign or size of \eqref{eq:fcrt-signed-weight} nor the location of
the collision in the compatible fibre.

Character orthogonality makes the missing analytic input precise.  Let
$\widehat G_S=\operatorname{Hom}(G_S,S^1)$.  For $g\in G_S$ and
$t\in\mathbb R$,
\[
 \sum_{\Phi_S(T)=g}e^{t y(T)}
 =\frac1{|G_S|}\sum_{\chi\in\widehat G_S}
   \overline{\chi(g)}
   \prod_{q\in J}\left(1+e^{t y_q}\chi(g_{S,q})\right).
\tag{FCRT-30}\label{eq:fcrt-one-fibre-fourier}
\]
Formula \eqref{eq:fcrt-one-fibre-fourier} only enumerates one prescribed
source fibre.  The FCRT flag instead requires a joint, nested selection of
$T$ in the $S$-fibre and a nonempty proper $U\subsetneq T$ in the
$\{p\}$-fibre, with $p\notin S$.  A future Fourier argument would therefore
need a rigorously defined two-fibre enumerator, control of its boundary
terms, and a uniform nontrivial-character estimate strong enough to retain
pairs with the required logarithmic weights.  It would then still have to
select several pairs subject to source and sink disjointness.  No such joint
estimate is asserted here.  Neither character orthogonality, the Boolean
collision count, nor the same-fibre exchange proposition supplies it.

\subsection{Scope of the exact finite computation}

The frozen computation uses integer factorizations and rational
multiplicative residuals; floating logarithms occur only as display fields.
It exhausts the normalized domain
\[
 1\le a\le b,\qquad a+b=c,\qquad\gcd(a,b)=1,\qquad c\le3000,
\tag{FCRT-31}\label{eq:fcrt-scan-domain}
\]
which contains $1{,}368{,}094$ triples.  The easy strata
$I=\varnothing$, $X\le Y$, and $|I|=1$ reduce the block search to proved
closed formulas: all boundaries vanish when $I=\varnothing$; SCRT and FCRT
vanish when $X\le Y$ because the full block is saturated, although PBT need
not vanish; and when $|I|=1$ with $X>Y$, all three boundaries equal
$X-Y$.  Only six points require a direct block search.  The producer
enumerates every compatible saturated block, represents SCRT states by all
attainable disjoint source--sink unions, retains every disjoint block family
for FCRT, and solves the residual ownership problems by an exact capped-state
dynamic program.  The independent validator instead recursively enumerates
every disjoint block family and then every residual sink and reuse-token
owner word.  It does not import the producer and performs a deterministic
byte-for-byte replay.  The six directly enumerated points are
\[
 \begin{gathered}
 (1,224,225),\ (1,675,676),\ (343,625,968),\\
 (1,2303,2304),\ (125,2187,2312),\ (99,2401,2500).
 \end{gathered}
\]
The same validation also covers $154$ fully factored structured-family rows,
including the out-of-range regression point $(1,65024,65025)$, and $17$
detailed certificates.

Besides the three rows in \eqref{eq:fcrt-three-witness-table}, the forced
comparison
\[
 \begin{array}{c|c|c|c|c}
 P&e^\Delta&e^{B_{\rm FCRT}}&e^{B_{\rm SCRT}}&e^{B_{\rm PBT}}\\ \hline
 (343,625,968)&44/35&44/35&11/7&11/7
 \end{array}
\tag{FCRT-32}\label{eq:fcrt-968-witness}
\]
is reproduced exactly.  Five points in the bounded domain have
$B_{\rm SCRT}>\Delta$; FCRT reaches the scalar lower bound at
$(1,675,676)$, $(343,625,968)$, and $(99,2401,2500)$.  The two remaining
FCRT fragmentation points are
$(1,224,225)$ and $(1,2303,2304)$.  The frozen validation status is
\texttt{PASS}; sealed hashes for the exact output, producer, validator, and
validation log are stored with the computation artifacts.

These statements have finite scope.  They prove no uniform upper bound and
produce no complete-premise unbounded counterfamily.  In particular, the
$n=284$ point in $(2,15^n-2,15^n)$ is recorded only with the proved fact
$v_{31}(15^{284}-2)=2$; its $335$-digit external arm is not completely
factored, so no exact FCRT optimum is asserted there.

\subsection{Lean formalization and the remaining arithmetic boundary}

The ordinary proofs above precede the formalization.  The primary Lean
module has $70$ public declarations and a matching $70$-query axiom audit.
An independently written bridge module has $18$ public declarations and a
matching $18$-query audit.  All four files compile with warnings treated as
errors.  The union reported by \texttt{\#print axioms} is exactly
\texttt{propext}, \texttt{Classical.choice}, and \texttt{Quot.sound}; no
\texttt{sorry}, \texttt{admit}, \texttt{unsafe}, or
\texttt{native\_decide} escape occurs.

The formalized unconditional content includes the finite maximal-feasible
subset and cancellation lemmas used in the nontrivial part of
Proposition~\ref{prop:fcrt-free-collapse}; an explicit owner kernel in which
each residual sink and reuse token has at most one owner; the clipped
residual and once-charge mass inequality; the cap of reuse credit by both
block surplus and witness weight; the endpoint height bridge; and the
conditional implication of the separately supplied uniform gate to the
repository's unchanged \texttt{ABCConjecture}.  The independent bridge caps
raw incoming credit by source demand,
\[
                 \widehat c_i=\min\{x_i,c_i\},
 \qquad x_i-\widehat c_i=\max\{x_i-c_i,0\},
\]
constructs the aggregate certificate from the owner kernel, and proves that
the clipped boundary is preserved exactly.  Its endpoint conversion assumes
the decomposition
\[
 \operatorname{signedEndpointCoreDefect}(P)
 =\operatorname{sourceMass}-\operatorname{sinkMass}
\tag{FCRT-33}\label{eq:fcrt-formal-endpoint-premise}
\]
as an explicit parameter.

Several formal limits are mathematically substantive.  The Lean selection
and cancellation lemmas do not encode the full typed free-target
configuration or all cases of Proposition~\ref{prop:fcrt-free-collapse}.
The formal \texttt{ProperSubfaceFlag} is not yet constructed from concrete
endpoint unitary divisors, and the owner kernel itself stores neither block
source sets nor the residual-target and block-disjointness proofs; it also
does not identify its abstract witness weight with $y(U)$.  A bare endpoint
certificate is tautologically inhabited by the boundary
$\max\{h-r,0\}$, so its mere existence carries no flagged-CRT content.
Concrete arithmetic admissibility, the endpoint identity
\eqref{eq:fcrt-formal-endpoint-premise}, and uniform smallness carry that
content.

The additive commutative-group versions of the complement identity and
same-fibre signed exchange are formalized.  The multiplicative residue-unit
construction \eqref{eq:fcrt-residue-group}--\eqref{eq:fcrt-packet-label}, the
inverse system, the Boolean collision criterion, and the Fourier identity
\eqref{eq:fcrt-one-fibre-fourier} are not yet formalized.  The required joint
cross-fibre enumerator has not yet been promoted to a theorem.  The witness
files check the displayed factorizations,
congruences, cap inequalities, and logarithmic strict inequalities, but do
not encode the exhaustive optimizer that proves the three exact minima in
Proposition~\ref{prop:fcrt-three-witnesses}.  Most significantly, no Lean
term inhabits the uniform gate \eqref{eq:fcrt-uniform-gate}.  Therefore this
checkpoint proves neither the standard $abc$ conjecture nor its negation;
the cross-fibre weighted selection estimate remains an open obligation.
